\section{Related Work}
Most work in fair machine learning deals with fairness with respect to single (binary) sensitive attributes. 
Multi-attribute fair classification was recently the focus of \citet{kearns2017preventing}---with empirical follow-up \citep{kearns2018empirical}---and \citet{johnson2018multicalibration}.
Both papers define the notion of an identifiable class of subgroups, and then obtain fair classification algorithms that are provably as efficient as the underlying learning problem for this class of subgroups.
The main difference is the underlying metric; \citet{kearns2017preventing} use statistical parity whereas \citet{johnson2018multicalibration} focus on calibration. 
Building on the multi-accuracy framework of \citet{johnson2018multicalibration}, \citet{kim2018multiaccuracy} develop a new algorithm to achieve multi-group accuracy via a post-processing boosting procedure. 

The search of independent latent components that explain observed data has long been a focus on the probabilistic modeling community \cite{comon1994independent,hyvarinen2000independent,bach2002kernel}.
In light of the increased prevalence of neural networks models in many data domains, the machine learning community has renewed its interest in learned features that ``disentangle'' semantic factors of data variation.
The introduction of the $\beta$-VAE \citep{higgins2016beta}, as discussed in section \ref{sec:background}, motivated a number of subsequent studies  that examine why adding additional weight on the KL-divergence of the ELBO encourages disentangled representations \citep{alemi2018fixing, burgess2018understanding}.
\citet{chen2018isolating,kim2018disentangling} and \citet{esmaeili2018structured} argue that decomposing the ELBO and penalizing the total correlation increases disentanglement in the latent representations. 
\citet{locatello2018challenging} conduct extensive experiments comparing existing unsupervised disentanglement methods and metrics. 
They conclude pessimistically that learning disentangled representations requires inductive biases and possibly additional supervision, but identify fair machine learning as a potential application where additional supervision is available by way of sensitive attributes.

Our work is the first to consider multi-attribute fair representation learning, which we accomplish by using sensitive attributes as labels to induce a factorized structure in the aggregate latent code.
\citet{bose2018compositional} proposed a compositional fair representation of graph-structured data.
\citet{kingma2014semi} previously incorporated (partially-observed) label information into the VAE framework to perform semi-supervised classification.
Several recent VAE variants have incorporated label information into latent variable learning for image synthesis \cite{klys2018learning} and single-attribute fair representation learning \cite{song2018learning,botros2018hierarchical,moyer2018invariant}.
Designing invariant representations with non-variational objectives has also been explored, including in reversible models \cite{ardizzone2018analyzing,jacobsen2018excessive}.
